\documentclass[11pt]{article}
\input{style-pc}

\begin{document}

\entetesujet{Sujet bac 2011 -- Physique-Chimie -- Série C}

% ============================ CHIMIE ============================
\matiere{CHIMIE}{8 points}

\exo{1}{}
On mélange dans plusieurs ampoules $3{,}7$ g d'acide propanoïque ($\ch{CH_3-CH_2COOH}$) et $1{,}6$ g de méthanol ($\ch{CH_3-OH}$). On scelle les ampoules et on les place dans une étuve à $50\,^\circ$C. Au bout de 24 heures, on constate que la masse d'acide propanoïque après la réaction reste constante et égale à $1{,}23$ g par ampoule.

\begin{enumerate}
  \item \begin{enumerate}[label=\alph*.]
    \item Quelle réaction chimique a eu lieu dans les ampoules ?
    \item Donner ses caractéristiques.
  \end{enumerate}
  \item \begin{enumerate}[label=\alph*.]
    \item Écrire l'équation-bilan de cette réaction.
    \item Donner le nom du composé organique formé.
  \end{enumerate}
  \item Calculer la quantité de matière (nombre de moles) du composé organique formé à l'équilibre.
  \item \begin{enumerate}[label=\alph*.]
    \item Calculer le rendement de cette réaction.
    \item Comment pourrait-on obtenir le même résultat expérimental en moins de temps ?
  \end{enumerate}
\end{enumerate}
\noindent On donne en g$\cdot$mol$^{-1}$ : C = 12 ; O = 16 ; H = 1.

\exo{2}{}
Au cours d'une séance de travaux pratiques, le professeur demande à un élève de préparer une solution $S_0$ d'ions $\ch{Fe^{2+}}$ en partant d'une masse $m = 13{,}9$ g de sulfate de fer II hydraté ($\ch{FeSO_4,7H_2O}$) qu'il dissout dans l'eau pure pour obtenir $500$ cm$^3$ de solution.

\begin{enumerate}
  \item Calculer la concentration molaire théorique $C_0$ de la solution $S_0$ obtenue.
  \item Afin de vérifier le travail effectué, le professeur demande à un autre élève de déterminer la concentration de la solution obtenue par dosage à l'aide d'une solution de permanganate de potassium $(\ch{K^+} + \ch{MnO_4^-})$, de concentration molaire $0{,}04$ mol$\cdot$L$^{-1}$.
  \begin{enumerate}[label=\alph*.]
    \item Écrire l'équation-bilan de la réaction qui a lieu.
    \item Sachant que $11$ cm$^3$ de la solution de permanganate de potassium ont été nécessaires pour doser $20$ cm$^3$ de la solution d'ions $\ch{Fe^{2+}}$, déterminer la concentration molaire volumique $C$ de la solution d'ions $\ch{Fe^{2+}}$.
    \item En déduire l'incertitude relative sur la concentration $C_0$.
  \end{enumerate}
\end{enumerate}
\noindent On rappelle que les couples redox en présence sont : $\ch{Fe^{3+}/Fe^{2+}}$ et $\ch{MnO_4^-/Mn^{2+}}$. On donne en g$\cdot$mol$^{-1}$ : Fe = 56 ; S = 32 ; H = 1.

% ============================ PHYSIQUE ============================
\matiere{PHYSIQUE}{12 points}

\exo{1}{}
$AB$ est une tige rigide de masse négligeable, de milieu $O$, de longueur $AB = 2L = 80$ cm. $AB$ peut osciller dans le plan vertical autour d'un axe $(\Delta)$ horizontal et passant par le point $O$. En $A$, on a fixé un solide de masse $M$ et en $B$ un solide de masse $m$ (ces deux solides sont ponctuels). On donne : $M = 300$ g ; $m = 100$ g et $g = 10$ m/s$^2$.

\begin{center}
\begin{tikzpicture}[>=Stealth,scale=1]
  % tige verticale
  \draw[thick] (0,-2)--(0,2);
  % axe Delta horizontal
  \draw[thick] (-2.2,0)--(2.6,0);
  \node[right] at (0.15,-0.3){$(\Delta)$};
  % masses
  \fill (0,2) circle (2.5pt) node[above]{$B(m)$};
  \fill (0,-2) circle (2.5pt) node[below]{$A(M)$};
  \fill (0,0) circle (1.2pt) node[above right]{$O$};
\end{tikzpicture}
\end{center}

\begin{enumerate}
  \item \begin{enumerate}[label=\alph*.]
    \item Calculer le moment d'inertie du système ainsi constitué par rapport à l'axe $(\Delta)$.
    \item Donner la position $G$ du centre d'inertie du système.
    \item On écarte ce système d'une faible amplitude de la position d'équilibre et on l'abandonne sans vitesse initiale.
    \begin{enumerate}[label=c\textsubscript{\arabic*}.]
      \item Établir l'équation différentielle du pendule ainsi constitué.
      \item En déduire la période du mouvement. Faire l'application numérique.
    \end{enumerate}
  \end{enumerate}
  \item Le pendule est écarté de sa position d'équilibre d'un angle $\alpha = 60^\circ$ et abandonné sans vitesse initiale.
  \begin{enumerate}[label=\alph*.]
    \item En utilisant le théorème de l'énergie cinétique, calculer la vitesse angulaire du pendule au passage à la position d'équilibre.
    \item En déduire la vitesse linéaire de $A$ à cette position.
  \end{enumerate}
\end{enumerate}

\exo{2}{}
Dans le dispositif d'Young, la source $S$ émet une radiation lumineuse de longueur d'onde $\lambda$ qui éclaire les fentes $S_1$ et $S_2$ distantes de $a = 7{\cdot}10^{-1}$ mm. On observe le phénomène d'interférences sur un écran situé à une distance $D = 1$ m du plan des fentes.

\begin{enumerate}
  \item Comment appelle-t-on la zone où l'on observe ce phénomène ?
  \item Sur l'écran, le milieu de la 7\textsuperscript{e} frange brillante est situé à $x = 4{,}2$ mm du milieu de la frange centrale. Calculer :
  \begin{enumerate}[label=\alph*.]
    \item L'interfrange $i$.
    \item La longueur d'onde $\lambda$ de la radiation.
  \end{enumerate}
  \item La source $S$ émet maintenant deux radiations, l'une de longueur d'onde $\lambda_1 = 0{,}420$ µm et l'autre de longueur d'onde inconnue $\lambda_2$.
  \begin{enumerate}[label=\alph*.]
    \item Décrire le phénomène observé sur l'écran.
    \item Sur l'écran, le milieu de la 8\textsuperscript{e} frange brillante de la radiation $\lambda_1$ coïncide avec le milieu de la 7\textsuperscript{e} frange brillante de la radiation $\lambda_2$. Calculer $\lambda_2$.
    \item Calculer la distance entre deux coïncidences successives.
  \end{enumerate}
\end{enumerate}

\exo{3}{}
On se propose de déterminer la résistance $r$ et l'inductance $L$ d'une bobine. Pour cela, on monte en série un conducteur ohmique de résistance $R = 7$ $\Omega$ et la bobine.

\begin{center}
\begin{tikzpicture}[>=Stealth,scale=1]
  % bornes
  \fill (0,0) circle (1.2pt) node[left]{$A$};
  \fill (9,0) circle (1.2pt) node[right]{$B$};
  % fil gauche vers R
  \draw (0,0)--(2,0);
  % resistance R
  \draw (2,-0.22) rectangle (3.2,0.22);
  \node[above] at (2.6,0.25){$R$};
  % fil vers noeud
  \draw (3.2,0)--(4.6,0);
  \fill (4.6,0) circle (1pt);
  \node[above] at (4,0.1){$i$};
  \draw[->] (3.9,0.0)--(4.3,0.0);
  % bobine (coil)
  \draw[decorate,decoration={coil,aspect=0.6,segment length=5pt,amplitude=4pt}] (4.6,0)--(6.6,0);
  \node[above] at (5.6,0.3){$(L,r)$};
  % fil vers B
  \draw (6.6,0)--(9,0);
  % U1
  \draw[<->] (0,-0.9)--(4.6,-0.9) node[midway,below]{$U_1$};
  % U2
  \draw[<->] (4.6,-0.9)--(9,-0.9) node[midway,below]{$U_2$};
  % U
  \draw[<->] (0,-1.6)--(9,-1.6) node[midway,below]{$U$};
  % pointillé au noeud
  \draw[dashed] (4.6,0)--(4.6,-0.9);
\end{tikzpicture}
\end{center}

\noindent L'ensemble est alimenté par une tension sinusoïdale de fréquence $N = 50$ Hz et de valeur efficace $U = 24$ V. On mesure les tensions efficaces $U_1$ et $U_2$ respectivement aux bornes du conducteur ohmique et aux bornes de la bobine. On obtient : $U_1 = 8$ V et $U_2 = 19{,}6$ V.

\begin{enumerate}
  \item \begin{enumerate}[label=\alph*.]
    \item Donner les expressions et calculer les impédances $Z_1$ du conducteur ohmique, $Z_2$ de la bobine et $Z$ du circuit.
    \item En déduire $r$ et $L$.
  \end{enumerate}
  \item On ajoute en série dans le circuit précédent un condensateur de capacité $C$. Le circuit étant capacitif :
  \begin{enumerate}[label=\alph*.]
    \item Quelle doit être la valeur de $C$ pour que l'intensité efficace soit la même que dans la question 1, la tension n'étant pas modifiée ainsi que la fréquence ?
    \item Exprimer la phase $\varphi$ de la nouvelle tension instantanée en fonction de $L$, $\omega$, $R$ et $r$ et en déduire $\varphi$.
    \item Construire le diagramme de Fresnel correspondant.
  \end{enumerate}
\end{enumerate}

\end{document}
